\section{Lecture 2. Is Our Universe Part of a Multiverse? II}

\subsection*{Summary of Last Lecture}

The Standard Big Bang: Really describes only the aftermath of a bang, beggining with a hot dense uniform soup of particles filling an expanding space.

Cosmic Inflation: The prequel, describes how repulsive gravity - a consequence of negative pressure - could have driven a tiny patch of the early Universe into exponential expansion. The total energy would be very small or maybe zero, with the negative energy of the cosmic gravitational field canceling the energy of matter:
\begin{enumerate}
    \item Inflation can explain the large-scale uniformity of the Universe. (Cosmic microwave background (CMB) uniform to 1 part in $100 000$)
    \item Inflation can explain why $\Omega \equiv \rho/\rho_{\rm crit}=1$ was accurrate to $>15$ decimal places at $t=1$ seconds. Predicts $\Omega = 1$. Data $\Omega = 1.0010 \pm 0.0065$.
    \item Predicts small quantum fluctuations in the mass density, which can be seen today as ripples in the CMB. Predictions agree very well with data.
\end{enumerate}

Most inflationary models become eternal - the expansion overpowers the decay of the repulsive gravity material, so inflation never ends. An exponential growing and never-ending number of pocker universes are formed where decays occur. The expansion of the Universe is accelerating, indicating that space is filled with "dark energy", most simply described as vacuum energy. Vacuum energy in a quantum field theory is not surprising - field fluctuations, nonzero Higgs field 0 there are positive and negative constributions. But typical magnitudes are $~ 10^{120}$ times too large.

The landscape of String Theory: String theory predicts ~ $10^{500}$ long-lived, metastable "vacua", any one of which can act as the vacuum for a pocket Universe. Each would have its own value for the vacuum energy density, with values ranging from roughly $-10^{120}$ to $+10^{120}$ times the observed value. If the landscape has $10^{500}$ vacua, and a fraction $10^{-120}$ have small vacuum energy densities like our universe, then we expect about

\begin{equation*}
    10^{-120} \times 10^{500} = 10^{380}
\end{equation*}
vacua with low energy densities like ours. How could we explain why we are living in such a fantastically unusual type of vacuum? Consider, as an exmaple, the local density of matter in which we find outselves - is about $10^{30}$ times larger than the mean density of the visible Universe. Most of us would presumably accept this as a selection effect: life can evolve only in those rare regions of the Universe where the density of matter is unusually high. As early as 1987, Steve Weinberg pointed out that the vacuum energy density might be explained in the same way. Maybe the vacuum energy density \textbf{is} huge in most pocket universes. Nonetheless, we need to remember that vacuum energy causes the expansion od the Universe quickly collapses. If large and positivem the universe flies apart before galaxies can form. It is plausible, therefore, that life can arise only if the vacuum energy density is very near zero.

In 1998 Martel, Shapiro and Weinberg made a serious calculation of the effect of the vacuum energy density on galaxy formation. The found that to within a factor of order 5, they could "explain" why the vacuum energy density is as small as what we measure.

Some consider these anthropic arguments as ridiculous. The anthropic explanation (for anything) should be considered the explanasion of last resort. Until we actually understand the landscape, and the initiation of life, we can only give plausibility arguments for anthropic explanations. Hence the anthropic arguments only become attractive wehen the search for more deterministic explanations has failed, as so far is the case for the vacuum energy density.